\section{An explicit witness: the separation on a rational datum}\label{sec:witness}

\subsection{The witness datum}\label{sec:datum}

Two architectures -- extraction ($q=0$) and regeneration ($q=1$) -- share one review interval $[0,1]$. The phase state is $(q,x,s_1,s_2)$: a physical reserve stock $x$ (a bridge stock, for instance an adjustment fund) and two typed floors $s_1$ and $s_2$ normalised at zero (on the benchmark of Section~\ref{sec:benchmark}: a biomass margin above a limit reference point and a fleet income margin). The transition-safe set is $S_0=\{x\ge0,\ s_1\ge0,\ s_2\ge0\}$ and the destination set is $G=\{(1,x,s):x\ge0,\ s\ge0\}$, with destination maintainability witnessed by a hold policy. The reset adds the strictly positive gain $e=(\tfrac14,\tfrac14)$, placing successors in the interior of $G$. The rescue cost is $c=1$. Initial states are $X_0=\{(0,x,s):x\ge0,\ s\ge0\}$.

Four plans are available (Table~\ref{tab:actions}), written as piecewise-linear maps on the breakpoints $\{0,\tfrac12,1\}$ and monotone on each piece, so that every tube is the exact visited set. The disturbance class is \emph{action-indexed}: the adverse disturbance triggers the worst-case dip, of fixed depth $2$, in the active plan's characteristic coordinate ($s_1$ for \act{fast}, $s_2$ for \act{slow}); on the benchmark the depth $2$ is a quota pulse of $3/2$ plus an environmental shock of $1/2$. The worst-case tube of \act{fast} is therefore $[s_1-2,s_1]$ in $s_1$ with $s_2$ and $x$ constant, that of \act{slow} is symmetric, and that of \act{staged} is $[x-1,x]$ in $x$ with both floors nondecreasing. A shared-shock class, in which the environmental component strikes every plan, is treated in Section~\ref{sec:commonshock}.

\begin{table}[!htbp]
\centering\small
\caption{The plan menu of the witness datum from an initial state $(0,x,s)$, $x\ge0$, $s\ge0$. Trajectories are worst-case; $e=(\tfrac14,\tfrac14)$.}
\label{tab:actions}
\begin{tabularx}{\textwidth}{@{}lXl@{}}
\toprule
Plan & Within-interval trajectory (worst case) & Successor \\
\midrule
\act{no-switch} & state constant & $\{(0,x,s)\}$ -- misses $G$ \\
\act{fast} & $s_1(t)=s_1-4t$ on $[0,\tfrac12]$, $s_1-4(1-t)$ on $[\tfrac12,1]$; $s_2,x$ constant & $\{(1,x,s+e)\}$ \\
\act{slow} & $s_2(t)=s_2-4t$ on $[0,\tfrac12]$, $s_2-4(1-t)$ on $[\tfrac12,1]$; $s_1,x$ constant & $\{(1,x,s+e)\}$ \\
\act{staged} & $x(t)=x-t$ on $[0,1]$; floors grow linearly to $s+e$ & $\{(1,x-1,s+e)\}$ \\
\bottomrule
\end{tabularx}
\end{table}

\noindent\textit{Role of \act{no-switch}.} \act{no-switch} is the status-quo, non-transition baseline: its successor keeps $q=0$ and therefore misses the destination set $G$. It remains in the four-action menu so that this failure is explicit; \act{fast}, \act{slow} and \act{staged} are the candidate transition plans.

The datum is deliberately spare. Its margin-class plans are \act{fast} and \act{slow}, with $\mathcal D=\{(2,0),(0,2)\}$, so Theorem~\ref{thm:menu} predicts the gap $\{s_1+s_2\ge2\}\setminus(\{s_1\ge2\}\cup\{s_2\ge2\})$ among them; \act{staged} lies outside the margin class, because its admissibility is governed by the reserve rather than by the floors, and it is what splits the predicted gap into a rescuable and an impossible part.

\subsection{The witnessed separation}

\begin{theorem}[witnessed separation]\label{thm:witness}
On the witness datum, over $X_0$:
\begin{enumerate}[label=\textup{(\arabic*)},leftmargin=2em]
\item $\Vtyp=\{x\ge1\}\cup\{s_1\ge2\}\cup\{s_2\ge2\}$.
\item $\Vweak=\bigcap_{w\in\Wp}\mathcal V_w=\{x\ge1\}\cup\{s_1+s_2\ge2\}$.
\item $\Vphys=\Vend=X_0$.
\end{enumerate}
Define the \emph{discrepancy region} $\Q=\{s_1<2,\ s_2<2,\ s_1+s_2\ge2\}$ ($x$ unrestricted), the \emph{rescue set} $R=\Q\cap\{x\ge1\}$ and the \emph{impossibility region} $I=\Q\cap\{x<1\}$. Then:
\begin{enumerate}[label=\textup{(\arabic*)},leftmargin=2em,resume]
\item $\FP=\Vweak\setminus\Vtyp=I$. The acceptance gap is exactly the impossibility region; it is nonempty, with relative interior $\{0<x<1,\ 0<s_1<2,\ 0<s_2<2,\ s_1+s_2>2\}$ in the $q=0$ slice. The rescue set $R$ lies in $\Vtyp$, not in the gap.
\item Both hierarchy inclusions are strict: every point of $I$ lies in $\Vweak\setminus\Vtyp$, and $(x,s_1,s_2)=(\tfrac12,\tfrac1{10},\tfrac1{10})$ lies in $\Vphys\setminus\Vweak$.
\item \textup{(Per-weight plan disagreement.)} On the relative interior of $\Q$ ($0<s_2<2$), for weights with $w_1>0$ put $r=w_2/w_1$. Then \act{fast} is licensed exactly for $r\ge\rho_1$ and \act{slow} exactly for $r\le\rho_2$, where
\[
\rho_1=\frac{2-s_1}{s_2},\qquad \rho_2=\frac{s_1}{2-s_2},\qquad \rho_2\ge\rho_1\iff s_1+s_2\ge2 .
\]
Among the two principal plans, $r>\rho_2$ licenses \act{fast} but not \act{slow}, $r<\rho_1$ licenses \act{slow} but not \act{fast}, and both are licensed for $\rho_1\le r\le\rho_2$. The axis weights are covered separately: $(1,0)$ licenses \act{slow} but not \act{fast}, while $(0,1)$ licenses \act{fast} but not \act{slow}, among those two plans. On $I$ no single plan serves every weight: $\bigcap_wE_w(z)=\Etyp(z)=\varnothing$. On $R$, \act{staged} is typed-admissible and serves every weight.
\item \textup{(The rescue split.)} $R$ is typed-transformable, witnessed by \act{staged}: bridging at cost $c=1$ keeps both floors intact and lands in $G$. $I$ is aggregate-feasible for every cone weight yet admits no typed-admissible plan, with four exhibited violations: \act{fast} breaches the $s_1$ floor under its adverse disturbance, \act{slow} the $s_2$ floor, \act{staged} drives $x$ negative, and \act{no-switch} misses $G$.
\end{enumerate}
\end{theorem}

The proof is in \ref{app:witness}; the formulas for $\rho_1,\rho_2$ require $0<s_2<2$, and the ratio is used for $w_1>0$; the two coordinate-axis weights are classified directly in part (6). On boundary faces in the state coordinates, the classification follows from the closed-set identities (1)--(2). Table~\ref{tab:regions} records the two verdicts region by region, and Figures~\ref{fig:triangle} and \ref{fig:paths} display them.

\paragraph{The witness point} At $(x,s_1,s_2)=(\tfrac12,\tfrac65,\tfrac65)$, the aggregate passes ($\tfrac{12}5\ge2$) but no menu plan is typed-admissible. The thresholds $(\rho_1,\rho_2)=(\tfrac23,\tfrac32)$ give \act{slow}-only, both, and \act{fast}-only licensing at $r=\tfrac12,1,2$, respectively. At $x=1$, \act{staged} rescues the same floors.

\paragraph{Reading} The compensatory assessment's binding condition is a total budget, $s_1+s_2\ge2$; the noncompensatory assessment requires either one floor to survive its own worst-case dip ($s_i\ge2$) or the reserve to cover the rescue cost ($x\ge1$). Between the two conditions the compensatory doctrine certifies -- per weight, with weight-dependent plans -- transitions the noncompensatory doctrine rejects, and the aggregate reading cannot tell $I$ from $R$: the two differ only in $x$, which the index does not see. Part (7) is a \emph{negative certificate} in the sense of a rejection with an exhibited violated constraint for each plan, exhausting the menu; it is a complete verdict relative to the menu, not an inconclusive one.

\begin{table}[!htbp]
\centering\small
\caption{Region dictionary on the witness datum (Theorem~\ref{thm:witness}).}
\label{tab:regions}
\begin{tabularx}{\textwidth}{@{}>{\raggedright\arraybackslash}p{0.30\textwidth}>{\raggedright\arraybackslash}p{0.17\textwidth}>{\raggedright\arraybackslash}p{0.14\textwidth}X@{}}
\toprule
Region & Per-weight verdict & Common-plan verdict & Reason \\
\midrule
$x\ge1$, all $s$ & accept & accept & \act{staged} serves every weight \\
$s_1\ge2$ or $s_2\ge2$ & accept & accept & \act{fast}/\act{slow} survives its own dip \\
$R$: $x\ge1$, $s_1<2$, $s_2<2$, $s_1+s_2\ge2$ & accept & accept & the reserve finances \act{staged} \\
$I$: $x<1$, $s_1<2$, $s_2<2$, $s_1+s_2\ge2$ & accept (weight-dependent plans) & reject & dips breach the floors; the reserve cannot finance phasing \\
$x<1$, $s_1+s_2<2$ & reject & reject & even the aggregate rejects \\
\bottomrule
\end{tabularx}
\end{table}

\begin{figure}[!htbp]
\centering
\includegraphics[width=\linewidth]{fig_witness_triangle.png}
\caption{The discrepancy triangle $\Q=\{s_1<2,\ s_2<2,\ s_1+s_2\ge2\}$ in the floor plane. Panel A ($x<1$): $\Q$ is the impossibility region $I=\FP$; the arrows show the worst-case dips of \act{fast} ($s_1\to s_1-2$) and \act{slow} ($s_2\to s_2-2$) crossing the floors from the witness point $(6/5,6/5)$, and the point $(1/10,1/10)$ is the second strictness witness. Panel B ($x\ge1$): the same triangle is the rescue set $R$, served by \act{staged}. Below the diagonal both assessments reject when $x<1$ and accept when $x\ge1$. The reserve bars show the bridge stock against the rescue cost $c=1$ at representative slices ($x=0$ and $x=1$; the witness point sits at $x=\tfrac12$). The bottom strip gives the licensing zones at the witness point: \act{slow} only for $r<\rho_1=\tfrac23$, both for $\tfrac23\le r\le\tfrac32$, \act{fast} only for $r>\rho_2=\tfrac32$.}
\label{fig:triangle}
\end{figure}

\begin{figure}[!htbp]
\centering
\includegraphics[width=0.9\linewidth]{fig_path_view.png}
\caption{Path view at the witness point under the worst-case disturbance. Rows: \act{fast}, \act{slow}, \act{staged}. Left: the floor paths (solid) against the zero floor (dashed). Right: the scalarized index $s_1+rs_2$ for $r\in\{2,1,\tfrac12\}$. Under \act{fast} the index stays nonnegative for $r=2$ and $r=1$ while $s_1$ breaches its floor; under \act{slow} the mirror holds. The only plan that protects both floors, \act{staged}, consumes the reserve linearly: infeasible at $x=\tfrac12$, feasible at $x=1$ -- the difference between $I$ and $R$.}
\label{fig:paths}
\end{figure}

\clearpage
\subsection{Persistence under horizon extension, and its limit}\label{sec:persist}

Is the one-interval separation an artefact of the one-interval framing? Two results delimit the answer.

\begin{proposition}[persistence and erasure]\label{prop:persist}
\textup{(i)} Extend the witness datum to $m\ge2$ review intervals by prepending hold intervals, with sole prefix action \act{hold} (tube $\{z\}$, successor $\{z\}$, safe set $S_0$) available at every prefix state, unchanged terminal and successor architecture, and no additional reset or disturbance branches. Then for every stage $j$, $\mathcal V^{\mathrm{typ}}_j\subseteq\bigcap_{w}\mathcal V^{w}_j\subseteq\mathcal V^{\mathrm{phys}}_j$, the stage-0 accepted regions are the witness regions pulled back through the holds, and both strictness witnesses of Theorem~\ref{thm:witness}(5) persist.
\textup{(ii)} Persistence is architecture-conditional: there is a two-interval datum whose single-interval truncation exhibits an acceptance gap and whose two-stage extension erases it. The same final typed and scalarized terminal rules are used in both versions; in the extension, a repair stage enlarges the stage-1 successor accepted sets to the whole state space, while the stage-1 tubes themselves remain safe.
\end{proposition}

The erasure datum of (ii) (SI~S2.2) has two floors, no reserve, and a first stage with two actions that transfer one unit from one floor to the other. At $(2/5,2/5)$, the one-stage version is weak-accepted for every weight (the actions are licensed for $r\ge3/7$ and $r\le7/3$) but typed-rejected. In the two-stage version the first-stage tubes stay in the nonnegative orthant, and a second-stage repair sends every successor to $(1,1)$, so both stage-1 accepted sets are all of $\mathbb R^2$; the original final typed and scalarized terminal rules are unchanged. The witness datum itself is not erased this way: on $I$ its plans violate the stage-local path constraint, which a later action cannot undo. The example shows one sufficient repair-stage mechanism, not a general erasure rule: safe continuation cannot repair an earlier tube breach, and equality of terminal sets or action coupling alone does not imply alignment of the pathwise constraints.

\subsection{What closes the gap and what does not}\label{sec:closing}

Under the witness's declared linear-flow extension, the physical convexification instrument is the \emph{closed blend family}: for $\delta\in[0,1]$, $\act{blend}_\delta$ uses the fractional control $u_\delta=\delta\,u_{\act{fast}}+(1-\delta)\,u_{\act{slow}}$. Its endpoints are the primitive plans, and its worst-case tube and successor are the corresponding convex combinations under the stated linearity assumption. This physical blend is stronger than abstract convexification of required-margin vectors; an application must demonstrate it from model equations or declare it as a management assumption (for example, fractional effort or quota allocation with monitoring and enforcement).

\begin{theorem}[blend collapse]\label{thm:blend}
On the witness datum, over $X_0$:
\begin{enumerate}[label=\textup{(\roman*)},leftmargin=2em]
\item $\act{blend}_\delta$ is typed-admissible at $z$ exactly when $\delta\in\big[\,1-\tfrac{s_2}2,\ \tfrac{s_1}2\,\big]\cap[0,1]$, an interval nonempty exactly when $s_1+s_2\ge2$ (a singleton at equality), independent of $x\ge0$ and of the weight;
\item the typed accepted set of the menu augmented by the blend family is exactly the compensatory region,
$\Vtyp^{\mathrm{blend}}=\{x\ge1\}\cup\{s_1+s_2\ge2\}=\Vweak$: the gap vanishes, and the blend family reaches nothing beyond $\Vweak$;
\item the admissible window depends on the state through $(s_1,s_2)$ alone, so a single $\delta$ serves every weight: on $I$ the blend is the common plan that no member of the finite menu supplies.
\end{enumerate}
\end{theorem}

\begin{proposition}[complete-excursion concatenation preserves primitive floor violations]\label{prop:timeshare}
Consider the control extension in which a sequential policy is required to execute each selected primitive at full strength for its complete unit-interval excursion, with visited set exactly the union of the corresponding primitive worst-case tubes. Within this specified class, a policy that uses both \act{fast} and \act{slow} is typed-admissible if and only if $s_1\ge2$ and $s_2\ge2$. Thus no such complete-excursion concatenation is typed-admissible on $I$.
\end{proposition}

This proposition is conditional on the stipulated union-of-complete-tubes semantics; it is not a theorem about ordinary finite-rate switching. For example, in the no-shock branch of a natural full-rate derivative extension $\dot s_1=-4$ before $t=1/2$ and $+4$ after for \act{fast}, with the analogous rates on $s_2$ for \act{slow}, the schedule fast on $[0,1/4]$, slow on $[1/4,1/2]$, fast on $[1/2,3/4]$, and slow on $[3/4,1]$ takes $(6/5,6/5)$ through $(1/5,6/5)$, $(1/5,1/5)$, $(6/5,1/5)$, and back to $(6/5,6/5)$. Both floors remain positive with three switches. This is a different admissible-control extension, not a counterexample to the conditional proposition.

The two results compare distinct declared control extensions. Under action-indexed disturbances and linear fractional flows, the blend family closes the typed gap exactly on $s_1+s_2\ge2$. Under the shared-shock class of Section~\ref{sec:commonshock}, the blend window instead requires $s_1+s_2\ge5/2$, so it is empty at the canonical point $s_1+s_2=12/5$ (Proposition~\ref{prop:commonshock}). The complete-excursion result concerns only a policy whose visited set is the union of both full primitive tubes; the three-switch example above demonstrates a different implementation. Neither statement identifies physical blending with a lottery over plans or ranks all admissible switching policies.

The remaining instrument acts on the reserve. Define the augmentation map $\mathsf{Aug}_\kappa:(x,s,\mathcal A)\mapsto(x+\kappa,\ s,\ \mathcal A\cup\{\act{staged}_\kappa\})$, where $\act{staged}_\kappa$ is \act{staged} run from the augmented reserve, and the \emph{rescue threshold} $\kappa^*(z)=\inf\{\kappa\ge0:\ \exists a\in\mathcal A_\kappa,\ a\in\Etyp(z)\}$. The augmented state is $(x+\kappa,s_1,s_2)$, so increasing $\kappa$ changes only the available reserve.

\begin{proposition}[rescue threshold]\label{prop:rescue}
On the witness datum, for every $z\in X_0$,
\[
\kappa^*(z)=(1-x)_+\cdot\mathbf 1[\,x<1,\ s_1<2,\ s_2<2\,]=(1-x)\cdot\mathbf 1[\,z\notin\Vtyp\,].
\]
\end{proposition}

On $I$ the threshold is $1-x>0$; on $R$ and on all of $\Vtyp$ it is $0$. Because $\act{staged}_\kappa$ is admissible exactly when $x+\kappa\ge1$, independently of the floors, the reserve route admits every state of the failure set $\{x<1,s_1<2,s_2<2\}$ at cost $1-x$ -- including states such as $(\tfrac12,\tfrac1{10},\tfrac1{10})$ that even the aggregate rejects; $I$ is the part of the failure set on which the aggregate still certifies. Along this route the distance to the feasible slice $\{x\ge1\}$ equals the violation $(1-x)_+$ itself, so the error-bound modulus \citep{Fabian2010} of the relation is identically $1$: the route is perfectly conditioned. Table~\ref{tab:instruments} summarises how the instruments act on the gap.

\begin{table}[!htbp]
\centering\small
\caption{Instruments acting on the acceptance gap.}
\label{tab:instruments}
\begin{tabularx}{\textwidth}{@{}>{\raggedright\arraybackslash}p{0.24\textwidth}>{\raggedright\arraybackslash}p{0.32\textwidth}X@{}}
\toprule
Instrument & Operation & Effect on the gap \\
\midrule
Reweighting & change $w$ & none: the gap is a statement over all $w$ \\
Restricting the weight family & replace $\Wp$ by $W\subset\Wp$ & $\mathcal V_W$ can only expand while $\Vtyp$ is unchanged, so $\mathcal V_W\setminus\Vtyp$ can only expand (Proposition~\ref{prop:monotone}) \\
Reserve augmentation & add $\act{staged}_\kappa$, $x+\kappa\ge1$ & closes it where only the reserve binds, at cost $\kappa^*=1-x$ (Proposition~\ref{prop:rescue}) \\
Menu convexification & replace $\mathcal D$ by $\conv(\mathcal D)$ (blend family) & closes it exactly on the scalarized region (Theorems~\ref{thm:menu}, \ref{thm:blend}) \\
Discrete time-sharing & union of the primitive tubes & closes it only where both dips are separately survivable (Proposition~\ref{prop:timeshare}) \\
\bottomrule
\end{tabularx}
\end{table}

\subsection{Beyond linear aggregation: the substitutability ladder}\label{sec:ladder}

The results so far concern the linear index. We evaluate Protocol~2 on the witness under a normalized-margin CES family. Set $\lambda_i=1+s_i$ and, for normalized weights $w_i\ge0$, $\sum_iw_i=1$, define
\[
\begin{gathered}
M_\theta(\lambda;w)=\Big(\sum_iw_i\lambda_i^\theta\Big)^{1/\theta},\qquad \sigma=\frac1{1-\theta},\\
-\infty<\theta<1\quad(0\le\sigma<\infty),\qquad \theta=1\text{ corresponds to }\sigma=\infty.
\end{gathered}
\]
with the usual limits: $\theta=1$ is the linear aggregate, $\theta=0$ the geometric mean, $\theta=-1$ the harmonic mean, and $\theta\to-\infty$ the weighted support minimum. The inputs $\lambda_i$ are constructed normalized-margin indices; they are not automatically physical abundance ratios. A state is accepted at rung $\theta$ when every weight has some menu plan satisfying $M_\theta\ge1$ along its worst-case tube, subject to the same physical and destination constraints. We impose an additional floor-safety convention: for every $\theta<1$, a plan is rejected if any tube coordinate has $\lambda_i\le0$, regardless of the weight. This is not an inherent property of every CES mean: for $0<\theta<1$, a zero input can still yield a positive mean, while negative inputs are outside the standard geometric-mean domain. The convention makes any collapsed floor coordinate independently disqualifying; the linear member alone is allowed to compensate across it. At $\theta=1$, $M_1\ge1$ is exactly $w\cdot s\ge0$, so the original compensatory protocol and its licensing thresholds are recovered. At the Leontief limit, $M_{-\infty}$ is the minimum over the positive-weight support. A strictly positive full-support weight therefore requires every coordinate to meet its floor along the tube, while typed safety passes every support minimum; hence the all-weight acceptance set is exactly the typed set.

On the gap-region diagonal $z=(x<1,s,s)$, $s\in(1,2)$, the per-weight cover condition reduces, by the symmetry of the two principal plans, to a single comparison at the equal weight: accept iff $(s-1)^\theta+(s+1)^\theta\ge2$ for $\theta>0$, iff $s^2\ge2$ for $\theta=0$, and iff $(s-1)^\theta+(s+1)^\theta\le2$ for $\theta<0$. Each rung therefore has an exact critical floor $s^*(\theta)$, and the critical floors form a ladder (Table~\ref{tab:ladder}, Figure~\ref{fig:ladder}).

\begin{theorem}[the substitutability ladder]\label{thm:ladder}
On the witness datum:
\begin{enumerate}[label=\textup{(\roman*)},leftmargin=2em]
\item \textup{(Nesting.)} For rungs $\theta>\theta'$ and every state, acceptance at $\theta'$ implies acceptance at $\theta$. The accepted rungs at each state form a prefix of the ladder, so the \emph{critical elasticity} $\sigma^*(z)=\inf\{\sigma:\ z\text{ accepted at }\sigma\}$ is well defined; it is strictly positive and finite on every gap state in the deep interior ($1<s_1<2$ and $1<s_2<2$), where the finite primitive menu has positive tube ratios and the scalarized margin is strict at every normalised weight; it is equal to $\infty$ on gap states with $\min(s_1,s_2)\le1$, for which only the linear member certifies. Boundary equalities are not covered by the strictness assertion.
\item \textup{(Pointwise.)} For every gap state in the strict interior of the floor-margin region ($0<s_1<2$, $0<s_2<2$, $s_1+s_2>2$), some strictly positive elasticity rejects it, so the Leontief limit is not needed to reject a fixed interior gap state. On the deep-interior subset ($1<s_1<2$ and $1<s_2<2$), a finite rung below the linear member still accepts the state. Gap states with $\min(s_1,s_2)\le1$ are accepted at the linear member only.
\item \textup{(Uniform on the diagonal slice.)} For every $\theta<1$ the rung falsely certifies an exact interval $[s^*(\theta),2)$ of the witness-defined diagonal gap states, with $s^*(\theta)<2$ strictly; for the displayed rungs the critical floors increase as elasticity falls: $1<\sqrt{6\sqrt3-9}<5/4<\sqrt2<\varphi<\sqrt3<\cdots\to2$ (Table~\ref{tab:ladder}), approaching the typed boundary. Thus no positive elasticity is uniformly safe even over the full gap (the diagonal intervals alone witness failure), $\inf_{z\in\Q}\sigma^*(z)=0$, and the only uniformly safe member is $\sigma=0$ -- the Leontief member, which is exactly the typed operator ($\mathcal V^{0}=\Vtyp$).
\end{enumerate}
\end{theorem}

\begin{table}[!htbp]
\centering\small
\caption{The exact ladder on diagonal gap states $z=(x<1,s,s)$: acceptance condition and critical floor $s^*(\theta)$ per rung.}
\label{tab:ladder}
\begin{tabularx}{\textwidth}{@{}ll>{\raggedright\arraybackslash}p{0.4\textwidth}X@{}}
\toprule
$\theta$ & $\sigma$ & accept iff & critical floor $s^*(\theta)$ \\
\midrule
$1$ & $\infty$ & $s\ge1$ & $1$ \\
$2/3$ & $3$ & $s^2\ge3$, else $27(s^2-1)^2\ge(3-s^2)^3$ & root in $(117/100,\,59/50]$ \\
$1/2$ & $2$ & $s\ge5/4$ & $5/4$ \\
$0$ & $1$ & $s^2\ge2$ & $\sqrt2$ \\
$-1$ & $1/2$ & $s^2\ge s+1$ & $\varphi=(1+\sqrt5)/2$ ($8/5$ rejects, $13/8$ accepts) \\
$-2$ & $1/3$ & $s^2\ge3$ & $\sqrt3$ \\
$-m$ & $1/(m+1)$ & $(s-1)^m+(s+1)^m\le2(s^2-1)^m$ & increasing in $m$ \\
$-\infty$ & $0$ & $s\ge2$ & $2$: the typed boundary \\
\bottomrule
\end{tabularx}
\end{table}

\begin{figure}[!htbp]
\centering
\includegraphics[width=\linewidth]{fig_substitutability_ladder.png}
\caption{The substitutability ladder (Theorem~\ref{thm:ladder}). (a) The critical floor $s^*(\theta)$ on the diagonal as a function of the elasticity $\sigma$; the shaded region above the curve is the set of diagonal gap states that the $\sigma$-aggregator falsely certifies. The exact algebraic rungs $1,\,5/4,\,\sqrt2,\,\varphi,\,\sqrt3$ are marked; the curve approaches the typed boundary $s=2$ only as $\sigma\to0$. (b) The critical elasticity $\sigma^*(s)$ along the diagonal (root of the master comparison), with the exact rational brackets of the deposited verifier overlaid: at the witness point $s=6/5$, $\sigma^*\in(2,3)$; at $s=5/4$, $\sigma^*=2$ exactly; $\sigma^*\to\infty$ as $s\to1^+$ and $\sigma^*\to0$ as $s\to2^-$.}
\label{fig:ladder}
\end{figure}

The theorem settles the objection in two scopes. Pointwise, the separation is not a Leontief phenomenon: at the witness point every aggregator with $\sigma\le2$ already rejects, and the floor-referenced dashboard value $(\lambda_1+\lambda_2)/2=(1/5+11/5)/2=6/5\ge1$ at the trough is contradicted by the geometric index $\sqrt{11/25}<1$. Here $\lambda_i=1+s_i$, so $(\lambda_1+\lambda_2)/2=1+(s_1+s_2)/2$ is tested against $1$. Section~\ref{sec:benchmark} reports the same verdict in floor-margin coordinates as the unnormalised sum $s_1+s_2$ at threshold $0$: for normalized equal weights $w=(1/2,1/2)$, this sum is twice $w\cdot s$, so the rescaling does not change acceptance. Uniformly on the witness-defined diagonal gap slice, for every positive elasticity there are gap states arbitrarily close to the typed boundary that the aggregator certifies and no plan completes safely. The appearance of $\sqrt2$, $\varphi$ and $\sqrt3$ as the exact critical floors of the geometric, harmonic and $\sigma=1/3$ members is, to our knowledge, a new exact witness for the substitutability debate. Off the diagonal, the geometric-member comparison has the exact criterion in SI~S3.2, with boundary $s_2=\rho(s_1)$. The finite-grid and numerical fixed-sum explorations in SI~S3 do not establish a global slice-minimum or stationary-point theorem, and no such global claim is used here. The identification of the $\theta=0$ member with the Living Planet Index is at the level of functional form only: published LPI inputs and weighting are not Protocol~2, and the constructed inputs $\lambda_i=1+s_i$ depend on the chosen positive normalization of each margin. Under the optional biomass scale map in Section~\ref{sec:benchmark}, $B_m/B_{\mathrm{lim}}=1+s_1/2$, not $1+s_1$; the index is therefore not a physical abundance ratio. The all-weights protocol is also more restrictive than acceptance at one specified published weight, so its false-acceptance exposure can be smaller.

\subsection{Beyond action-indexed disturbances: the common shock}\label{sec:commonshock}

A second objection is that the gap depends on the action-indexed disturbance convention, under which the adverse event strikes only the coordinate the active plan moves. Decompose the datum as on the benchmark: each principal plan's announced schedule is a mid-interval tent of depth $3/2$ in its characteristic coordinate, and each environmental event is a transient tent dip of depth $\delta$ in one floor coordinate (the datum's depth $2=3/2+1/2$ fixes $\delta=1/2$); \act{staged}'s floors grow linearly, so an event of depth $\delta$ dips a growing floor by $\max(0,\delta-1/8)$. Under the \emph{common-shock class} the two coordinate events are separate disturbances each of which strikes every plan -- a marine heatwave striking the same biomass coordinate regardless of which harvest plan is running.

\begin{proposition}[the common-shock verdict]\label{prop:commonshock}
Under the common-shock class at $\delta=1/2$:
\begin{enumerate}[label=\textup{(\roman*)},leftmargin=2em]
\item $\Vtyp^{\mathrm{cs}}=\{s_1\ge2,\ s_2\ge\tfrac12\}\cup\{s_2\ge2,\ s_1\ge\tfrac12\}\cup\{x\ge1,\ s_1\ge\tfrac38,\ s_2\ge\tfrac38\}$ and
$\Vweak^{\mathrm{cs}}=\{x<1,\ s_1\ge\tfrac12,\ s_2\ge\tfrac12,\ s_1+s_2\ge2\}\cup\{x\ge1,\ s_1\ge\tfrac38,\ s_2\ge\tfrac38\}$, so the gap
$\mathrm{FP}_{\mathrm{cs}}=\{x<1,\ \tfrac12\le s_1<2,\ \tfrac12\le s_2<2,\ s_1+s_2\ge2\}$
survives: it is the action-indexed gap $I$ reduced by the exact $1/2$-cut, nonempty with interior, and the witness point is a member. On the surviving region the licensing thresholds are unshifted.
\item The rescue threshold becomes a trichotomy: $\kappa^*_{\mathrm{cs}}=0$ on $\Vtyp^{\mathrm{cs}}$; $(1-x)_+$ on failure states with $\min(s_1,s_2)\ge\tfrac38$; and $+\infty$ on failure states with $\min(s_1,s_2)<\tfrac38$, which no reserve increment rescues (witness $(\tfrac12,\tfrac14,3)$).
\item The blend window is nonempty exactly on $\{s_1\ge\tfrac12,\ s_2\ge\tfrac12,\ s_1+s_2\ge\tfrac52\}$, so the common-shock gap is not blend-closed on the band $s_1+s_2\in[2,\tfrac52)$ (witness $(\tfrac12,1,1)$).
\end{enumerate}
\end{proposition}

The separation is not an artefact of the action-indexed convention: under the common-shock model the gap persists, reduced by the exact $1/2$-cut. The model's rescue set also changes: on failure states with $\min(s)<3/8$, no finite reserve increment certifies, while the staged route has the stated coordinate-floor requirement. These are conditional feasibility results for the specified menu and disturbance class, not operational prescriptions. The coupled class (one disturbance carrying both events), its depth dependence with exact boundary $\delta=5/4$, the relocation of the gap to larger floors at full magnitude, and the collapse of the separation under a replace reading in which every plan faces the same plan- and weight-independent dip are stated and proved in SI~S4; the mechanism of that replace-reading collapse is that a common plan- and weight-independent dip removes the weak test's advantage of weight-wise plan choice, so the per-weight cover reduces to the typed condition.

\subsection{Verification}\label{sec:verification}

The retained analytic results are proved in Appendix~\ref{app:general}, Appendix~\ref{app:witness}, and SI~S2--S4. Exact-integer and rational verifiers provide finite supporting checks only; their state grids and check lists do not establish continuum claims, and the availability of each script and log is stated in SI~S7 and Data Availability. The cited Lean module contains formalized algebraic witness components and selected CES rungs, but its own scope excludes the general real-exponent nesting, continuity/IVT transfer, and relative-interior argument used by the full continuum theorem. We therefore treat those steps as paper proofs, not as machine-checked Lean results. No Lean build was run for this v77 revision; SI~S7 records the claim-to-coverage limits and the version mismatch of the available deposits.
